\chapter{Literature Review}
\label{ch:literature}

\section{Introduction}

This chapter surveys the bodies of work that directly shaped the research. Chinese smart park implementations provided the design model. Offline-first and PWA literature supplied the technical approach. ICT4D and e-government scholarship grounded the development context. Design Science Research gave the methodology its structure. Each domain is reviewed in turn, and the chapter closes with a gap analysis showing what the existing literature leaves open.

\section{Chinese Smart Park Digital Ecosystems}

\subsection{Evolution of Smart Industrial Parks in China}

Chinese industrial parks have changed substantially over forty years. \citet{Yang2025} identify three phases: an establishment period through the 1990s that prioritised attracting foreign capital; an optimisation period in the 2000s and 2010s focused on industrial clustering; and a smart phase from around 2015 onward defined by digital transformation \citep{Yu2020, Zeng2021, Chen2022}.

The shift is not just technological. \citet{Huang2023} describe it as a change in how parks understand what they are for --- moving from ``space provider'' to ``service ecosystem.'' Parks that once delivered land and utilities now manage tenant operations through digital platforms.

In practice, this means layers of technology working together: IoT sensors for environmental monitoring, digital platforms for administrative services, analytics tools for operational management, and connections to national e-government systems \citep{Wang2022}. The architecture is not standardised across parks, but the basic logic is consistent.

\subsection{Digital One-Stop Service Platforms}

The one-stop service model is what most clearly distinguishes Chinese smart park platforms from earlier administrative arrangements. Rather than sending tenants to separate departments for different needs, parks built integrated interfaces that handle multiple services in one place --- an approach borrowed from broader e-government reforms and adapted for the industrial park setting \citep{Yang2023}.

\citet{Schreieck2022} examine platform strategies in Chinese industrial parks and identify three common approaches: proprietary platform development, adoption of commercial enterprise platforms (WeChat Work, DingTalk), and hybrid models combining aspects of both. Parks using commercial platforms benefit from established user familiarity and mature feature sets; proprietary platforms, meanwhile, offer greater customization.

WeChat Work (internationally known as WeCom) has emerged as a dominant vehicle for smart park digital services. As of 2021, \citet{Tencent2024} report over 5.5 million enterprise users, with industrial parks forming a notable segment. The platform provides enterprise messaging, workflow automation, approval routing, and integration with the broader WeChat ecosystem --- including payment services. DingTalk, Alibaba's alternative, offers comparable functionality with particular strengths in workflow management and cloud integration. \citet{Alibaba2024} document features such as OA approval workflows, collaborative document editing, attendance tracking, and video conferencing for up to 1,000 participants, with low-code tools enabling rapid customization.

\subsection{Park Innovation Management Framework}

\citet{Yang2025} propose the Park Innovation Management (PIM) framework as a lens for understanding digital transformation in Chinese industrial parks. Five interconnected dimensions structure the framework: institutional innovation (governance structures and policies), technological innovation (digital infrastructure and platforms), service innovation (new service delivery models), ecosystem innovation (stakeholder relationships and value networks), and spatial innovation (bridging the physical and digital).

For this research, the PIM framework is useful because it maps which elements of Chinese smart park models might transfer and which will require rethinking. More importantly, it shows that technology transfer cannot work by copying the software alone --- the institutional, cultural, and service innovation dimensions matter just as much.

\subsection{Suzhou Industrial Park Case Study}

Suzhou Industrial Park (SIP), launched in 1994 as a China-Singapore collaboration, represents one of the most mature smart park implementations in operation. Its ``Marvelous City of SIP'' mobile application functions as a comprehensive one-stop service platform, available in both English and Chinese \citep{PRNewswire2024}.

Services span business registration, permits, utilities, and cultural amenities. Key transferable principles include a single service access point, mobile-first design, integrated payment workflows, and transparent request tracking --- alongside a dependency on reliable connectivity and Chinese payment ecosystems that does not transfer directly.

\section{Belt and Road Initiative and Digital Silk Road}

\subsection{China-Africa Technology Transfer Context}

Launched in 2013, the Belt and Road Initiative (BRI) established extensive frameworks for infrastructure development and technology transfer between China and partner countries. The Digital Silk Road (DSR), formally announced in 2015, extended these frameworks specifically to digital infrastructure, e-commerce, and smart city technologies \citep{ElKadi2024}.

\citet{deKluiver2024} review BRI and DSR implications for African digital development. The picture is mixed: there are genuine opportunities --- access to Chinese technology, financing, and working governance models --- but also legitimate concerns around technological dependency and contextual fit. Ethiopian industrial parks are deep inside these China-Africa relationships. Chinese capital and expertise built several of them, Chinese companies remain among the largest tenants, and technical ties between Ethiopian and Chinese park managers are ongoing \citep{Negesa2022, Brautigam2011}. That relationship creates a practical basis --- and a reasonable motivation --- for asking whether Chinese digital governance tools could be adapted to work there.

\subsection{Technology Adaptation in South-South Transfer}

The South-South technology transfer literature is consistent on one point: copying does not work. \citet{Chen2018} argue that receiving countries need ``absorptive capacity'' --- the ability to recognise, process, and genuinely apply external knowledge rather than simply import it. \citet{Hensengerth2018}, studying China-Africa infrastructure transfers specifically, found that projects with localised training, adjusted designs, and continued support tended to succeed where wholesale imports did not. Digital transfers add another layer of difficulty: local digital ecosystems, payment infrastructure, and connectivity conditions all differ, and those differences are not cosmetic.

\section{Digital Governance in Developing Countries}

\subsection{E-Government Evolution and Models}

E-government research has moved through several phases. Early work by \citet{Layne2001} described a four-stage model --- cataloguing, transactions, vertical integration, horizontal integration --- that set the terms for much subsequent thinking. \citet{Wimmer2002} brought a European lens focused on cross-boundary integration. \citet{Bannister2015} tracks the field's more recent turn toward citizen-centred design and governance goals that go beyond efficiency.

The UN E-Government Survey places Ethiopia in the middle tier on its Development Index, noting better online service availability alongside persistent gaps in telecommunications infrastructure and digital skills \citep{UN2022}. That combination --- improving but still constrained --- describes the environment this research operates in.

\subsection{E-Government in Ethiopia}

Ethiopian e-government has expanded since the Ministry of Innovation and Technology was established, though unevenly. \citet{Lessa2019} found infrastructure constraints as the primary barrier, with mobile-based delivery providing a partial workaround. \citet{Denbu2019}, studying public organisations specifically, identified connectivity problems, limited digital literacy, and institutional resistance as the main obstacles, and recommended mobile-first, offline-capable designs as the pragmatic path forward --- exactly the design logic this research follows.

Ethiopia's Digital Strategy 2025 identifies industrial parks as priority zones for digital infrastructure investment. The national policy direction and this research's objectives point toward the same gap.

\subsection{One-Stop Service Delivery Models}

One-stop service models have been implemented widely in e-government. \citet{Askim2011} document consistent benefits across national cases: service consolidation, lower administrative burden, better coordination. \citet{Bhatnagar2014} found similar efficiency gains and higher satisfaction in South Asian developing country implementations, though inter-agency coordination remained a persistent challenge.

In an industrial park, the one-stop model addresses a very specific problem: tenants currently have to visit different departments for different services. A unified platform removes that burden. The coordination still has to happen, but it happens inside the system rather than falling on the tenant to manage.

\section{Offline-First Architecture and PWA Technologies}

\subsection{Local-First Software Principles}

\citet{Kleppmann2019} lay out the local-first software paradigm as a deliberate challenge to conventional web application design. Where standard applications treat the server as the authoritative data source, local-first software makes locally stored data authoritative and uses the network opportunistically for synchronization rather than as a prerequisite for operation.

Seven ideals frame the paradigm \citep{Kleppmann2019}: work happens locally without waiting for the network; data is not trapped on a single device; the network is optional; collaboration with colleagues is seamless; data remains accessible over the long term even if service providers disappear; security and privacy are defaults; and users retain ultimate ownership of their data. In environments where connectivity cannot be taken for granted, these principles have direct practical value --- people can get their work done regardless of network conditions.

\subsection{Progressive Web Application Technologies}

Progressive Web Applications combine a set of technologies and design patterns that allow web applications to deliver app-like experiences while preserving the reach and accessibility of the web \citep{AterPWA2017, Majchrzak2018}. \citet{Tandel2018} document measurable gains in user engagement and retention over traditional web applications, attributing these to three key capabilities: home screen installation, offline operation, and progressive enhancement --- functioning on any browser while offering richer features where supported.

Technically, PWAs rest on three pillars. Service workers are JavaScript files that run in a separate thread, enabling background processing, request interception, and caching. Web app manifests supply the metadata needed for installation and customization. Caching APIs --- particularly the Cache API and IndexedDB --- provide persistent client-side storage for assets and data alike.

\citet{BiornHansen2017} compare PWAs with native mobile applications across performance, capability, development cost, and user experience, finding that PWAs deliver comparable results for many application types at significantly lower development and deployment costs. The advantage is especially pronounced when offline capability and multi-platform reach are priorities.

\subsection{Service Workers and Caching Strategies}

Service workers enable the caching strategies that offline-first applications depend on \citep{Gaunt2020}. Google's Workbox library provides production-ready modules for managing service worker logic, simplifying the implementation of common patterns \citep{Workbox2024}. \citet{ArchibaldJake2023} catalog these patterns extensively in the ``Offline Cookbook,'' while \citet{AhnAllen2020} describe the principal approaches: cache-first (serve from cache, fall back to network), network-first (try network, fall back to cache), stale-while-revalidate (serve immediately from cache while updating in the background), and cache-only or network-only strategies for specific resource types.

For offline-first applications, cache-first is generally preferred for assets and previously loaded data, since it guarantees availability regardless of connectivity. Stale-while-revalidate suits data that changes periodically, balancing instant responsiveness with eventual consistency. Background synchronization, enabled through the Background Sync API, lets PWAs defer network operations until connectivity returns --- essential for queueing data submissions created while offline.

\subsection{IndexedDB and Client-Side Data Persistence}

IndexedDB provides a low-level API for storing substantial amounts of structured data on the client. Unlike localStorage --- limited to string key-value pairs and typically capped at 5--10\,MB --- IndexedDB supports complex data structures, indexing, and storage limits that typically reach hundreds of megabytes \citep{MDN2024}. In offline-first applications, it serves as the local database: pending changes accumulate there and synchronize with the server when connectivity resumes. Libraries such as Dexie.js \citep{Dexie2024} abstract over the raw IndexedDB API, simplifying common operations and improving code maintainability. \citet{Nicoara2018} offer detailed guidance on offline-first development patterns using these technologies.

\subsection{Firebase Offline Persistence}

Google's Firebase platform provides built-in offline persistence through Firestore. When enabled, Firestore automatically caches data locally and synchronizes changes once the network is available again \citep{GoogleFirebase2024}, offloading most synchronization logic to the platform rather than requiring application-level implementation. Conflicts are resolved using a last-write-wins strategy at the document level; more complex resolution can be implemented via Cloud Functions triggered by database changes. Firebase's generous free tier also makes it practical for thesis research and prototype deployment.

\section{ICT4D and Connectivity Constraints}

\subsection{ICT for Development Frameworks}

ICT4D research asks how digital technologies can contribute to development in resource-constrained settings \citep{Avgerou2010, Walsham2017}. \citet{Heeks2018} describes the field's progression from ICT4D 1.0 --- transferring technology from developed countries --- through participatory and locally adapted approaches, to a third wave that embeds technology questions within broader development processes \citep{Zheng2018, Sahay2017, Sein2019}. The one finding that holds across all phases: technology interventions need to fit local conditions. \citet{Toyama2015} puts it in terms of amplification --- technology magnifies what human and institutional capacity already exists, but it cannot create that capacity from nothing.

\subsection{Connectivity Challenges in Developing Regions}

Connectivity infrastructure in developing regions often diverges fundamentally from the assumptions embedded in technologies designed for reliably connected environments. \citet{ChenNath2016} document that intermittent access --- periods of availability interspersed with outages and variable quality --- is the norm, not the exception. Delay-tolerant networking research offers relevant technical frameworks. \citet{Fall2003} introduced delay-tolerant networking for scenarios where continuous end-to-end connectivity cannot be assumed; though originally targeting extreme settings like disaster response and interplanetary networks, the underlying principles apply directly to everyday connectivity challenges in developing regions.

\subsection{Mobile-First and Offline-First in Development Contexts}

The rapid spread of mobile devices in developing regions has driven sustained interest in mobile-first design. \citet{Donner2015} find that mobile devices frequently serve as people's primary --- and sometimes only --- computing platform. Offline-first design extends this mobile-first orientation to address connectivity limitations more directly. \citet{Muawwal2024} demonstrate that offline-first application design meaningfully improves both service accessibility and user satisfaction in government service contexts with limited connectivity.

\citet{Cherukuri2024} provide practical implementation guidance for offline-first PWAs, covering synchronization patterns, conflict resolution, and UX design for offline scenarios. A key recommendation: give users clear, continuous feedback about connectivity status and synchronization state.

\section{Technology Adaptation Frameworks}

Two theoretical threads inform the adaptation approach taken here. Appropriate Technology Theory, rooted in \citet{Schumacher1973}, holds that solutions must suit local conditions --- available resources, skill levels, and institutional realities --- rather than assuming that advanced technologies from industrialized settings are universally applicable. \citet{Toyama2015} extends this to digital contexts: technology amplifies existing human and institutional capacity; it cannot create what is not already there. Classical technology transfer theory reinforces the point: \citet{Rogers2003} identifies absorptive capacity and institutional environment as key determinants of success, while \citet{Edgerton2007} shows that adoption is never passive --- recipients always adapt technologies as they deploy them.

On the platform side, \citet{Schreieck2022} describe adaptation strategies along dimensions of feature modification, integration adaptation, and ecosystem repositioning. The adaptation in this research is substantial on all three: the connectivity architecture changes fundamentally, the payment ecosystem is replaced entirely with an internal token system, and core functionality (service consolidation, workflow automation, status tracking) is preserved while being reshaped for a different institutional context.

\section{Ethiopian Industrial Parks Context}

\subsection{Industrial Parks Development Corporation (IPDC)}

The Industrial Parks Development Corporation was established in 2014 as a public enterprise responsible for developing and managing Ethiopia's industrial parks. It is the institutional instrument behind the country's push to shift its economy from agriculture toward manufacturing \citep{Oqubay2018, Guteta2024}. \citet{Gebreeyesus2019} traces how this strategy developed, placing the park programme within a wider effort at structural economic transformation.

IPDC's portfolio includes Bole Lemi Industrial Park (the first to become operational, in 2014), Eastern Industrial Zone (developed with Chinese partners), Hawassa Industrial Park (the flagship eco-industrial facility), and several parks at various development stages. These parks host both Ethiopian and international manufacturers, with a notable presence of Chinese, Turkish, and Indian firms \citep{Tesfaw2023}.

\subsection{Current Service Delivery Challenges}

\citet{Guteta2022} document recurring shortfalls in administrative responsiveness, infrastructure reliability, and communication at Ethiopian industrial parks. \citet{Guteta2023} used SERVQUAL to measure the gap between what tenants expect and what they actually experience, finding shortfalls across every dimension and recommending digital systems as the path toward improvement.

The operational reality is straightforward: processes run on paper. Requesting a service means visiting an office. Approvals follow paper through manual routing. Checking on a request means making a phone call. Tenants who came from places where this is handled digitally --- China, India, Turkey --- find the contrast difficult to overlook, and research confirms it shapes how they assess IPDC.

\subsection{Infrastructure and Connectivity Context}

Ethiopian telecommunications infrastructure has grown substantially in recent years, with mobile penetration surpassing 50\% and internet usage rising quickly. Connectivity remains unreliable outside major cities, however, and industrial parks --- despite their strategic importance --- are not insulated from regular disruptions. Survey respondents consistently reported three to five hours of daily internet disruption in their parks, with complete outages occurring two to three times monthly. This profile is fundamentally incompatible with cloud-dependent systems that treat network availability as a given.

Government investment in telecommunications --- fiber optic expansion, 4G coverage --- is underway, but infrastructure improvements take time. Any near-term digital solution must work within current constraints rather than banking on connectivity improvements that have yet to materialize.

\subsection{China-Ethiopia Industrial Park Relationships}

China-Ethiopia ties in industrial park development span investment, construction, tenancy, and technical cooperation \citep{Oqubay2019}. \citet{Negesa2022} document the scale of Chinese involvement in financing and building Ethiopian parks. Chinese companies form a substantial share of tenants --- particularly in textiles and garments --- meaning park stakeholders already have direct familiarity with Chinese business practices and may be receptive to digital systems informed by Chinese models. Technical cooperation channels between Ethiopian and Chinese park managers provide further routes for knowledge transfer, while also sharpening attention to how much genuine adaptation is needed to suit Ethiopian conditions.

\section{Design Science Research Methodology}

\subsection{DSR Paradigm and Principles}

Design Science Research is built around the idea that creating and evaluating artifacts is itself a valid form of research \citep{Hevner2004}. Where behavioural research tries to understand and explain phenomena, DSR tries to solve problems through purposeful design. Seven guidelines govern the approach: the research must produce a working artifact, address a real and relevant problem, evaluate utility rigorously, make clear knowledge contributions, apply methods systematically, treat design as a search through possibilities, and communicate results to appropriate audiences.

\subsection{DSRM Process and Artifact Types}

\citet{Peffers2007} propose a six-activity DSRM process model --- problem identification, defining objectives, design and development, demonstration, evaluation, and communication --- that organizes this research from literature review through thesis. \citet{March1995} identify four artifact types (constructs, models, methods, instantiations); this research produces all three non-trivial types: an instantiation (the platform prototype), a method (the adaptation framework), and constructs (the design principles). \citet{Gregor2013} argues that design theory --- knowing how to do something --- is epistemically as valuable as explanatory theory, a position that justifies DSR as the appropriate paradigm where existing solutions fall demonstrably short.

\section{Related Works and Comparative Analysis}

\subsection{Smart Park Implementations}

Several research efforts have examined smart park implementations in various settings. Table~\ref{tab:related-smart-park} summarizes the most relevant works and their relationship to this research.

\begin{table}[H]
\centering
\caption{Related Smart Park Research}
\label{tab:related-smart-park}
\begin{tabular}{p{3.5cm}p{2.5cm}p{3.0cm}p{3.0cm}}
\toprule
\textbf{Study} & \textbf{Focus} & \textbf{Context} & \textbf{Gap Addressed} \\
\midrule
Yang et al. (2025) & Park Innovation Mgmt. framework & Chinese smart parks & Conceptual framework only \\
Huang et al. (2023) & SEZ digitalization & China & Assumes connectivity \\
Wang et al. (2022) & Smart park platforms & China & Assumes connectivity \\
Grosse-Bley \& Kostka (2021) & Big data in smart cities & Shenzhen, China & Does not address offline \\
\textbf{Current Research} & \textbf{Offline-first adaptation} & \textbf{Ethiopia} & \textbf{Addresses connectivity constraints} \\
\bottomrule
\end{tabular}
\end{table}

\subsection{Offline-First and PWA Applications}

Research on offline-first and PWA applications in developing country contexts offers useful technical guidance while revealing gaps this research aims to fill. Table~\ref{tab:related-pwa} summarizes the relevant PWA studies.

\begin{table}[H]
\centering
\caption{Related PWA and Offline-First Research}
\label{tab:related-pwa}
\begin{tabular}{p{3.5cm}p{2.5cm}p{3.0cm}p{3.0cm}}
\toprule
\textbf{Study} & \textbf{Focus} & \textbf{App.\ Domain} & \textbf{Gap} \\
\midrule
Cherukuri et al. (2024) & Offline-first PWA patterns & General & Not industrial context \\
Muawwal et al. (2024) & Offline gov.\ services & Citizen services & Simple workflows only \\
Ahn \& Allen (2020) & PWA capabilities & General web apps & Not developing country \\
\textbf{Current Research} & \textbf{Offline-first industrial OSS} & \textbf{Industrial parks} & \textbf{Complex workflows, Ethiopia} \\
\bottomrule
\end{tabular}
\end{table}

\subsection{E-Government in Ethiopian Context}

Ethiopian e-government research provides important context while making clear the absence of industrial park-specific digital solutions. \citet{Lessa2019} and \citet{Denbu2019} document broader e-government challenges in Ethiopia without addressing industrial park service delivery. This research extends that body of work into the industrial park domain, tackling the connectivity constraints that hold back conventional e-government approaches.

\section{Research Gap Analysis}

\subsection{Identified Gaps}

Four gaps emerge from the review. Table~\ref{tab:gap-matrix} maps them against the literature domains.

\textbf{Gap 1:} Chinese smart park implementations are well documented, but no adaptation framework exists for connectivity-constrained environments. The literature assumes persistent broadband.

\textbf{Gap 2:} Offline-first and PWA research covers simple single-user services. The multi-stakeholder workflows of industrial park service delivery --- approvals, status tracking, complaint resolution --- have not been addressed.

\textbf{Gap 3:} No purpose-built digital solution exists for Ethiopian industrial parks. E-government research covers Ethiopia but not industrial parks; industrial park research covers Ethiopia but not connectivity constraints.

\textbf{Gap 4:} General offline-first design patterns exist but industrial-specific guidance does not: which features to prioritize for offline use, how to sync multi-stakeholder data, how to present connectivity state to operational users.

Table~\ref{tab:gap-matrix} presents a gap matrix mapping how existing literature addresses --- or fails to address --- the key dimensions of this research.

\begin{table}[H]
\centering
\caption{Research Gap Matrix}
\label{tab:gap-matrix}
\begin{tabular}{lcccc}
\toprule
\textbf{Literature Domain} & \textbf{Smart Park} & \textbf{Offline-First} & \textbf{Ethiopia} & \textbf{Industrial} \\
\midrule
Chinese Smart Parks & \checkmark & $\times$ & $\times$ & \checkmark \\
PWA/Offline-First & $\times$ & \checkmark & $\times$ & $\times$ \\
Ethiopian E-Gov & $\times$ & $\times$ & \checkmark & $\times$ \\
ICT4D & $\times$ & Partial & Partial & $\times$ \\
\textbf{Current Research} & \checkmark & \checkmark & \checkmark & \checkmark \\
\bottomrule
\end{tabular}
\end{table}

\section{Conceptual Framework}

Figure~\ref{fig:conceptual-framework} maps the relationship between the bodies of literature reviewed in this chapter. Chinese smart park models supply the design vocabulary; technology adaptation theory governs how that vocabulary travels; offline-first architecture answers the connectivity constraint; and the Ethiopian IPDC context sets the requirements the resulting platform must satisfy.

\begin{figure}[H]
\centering
\resizebox{0.95\textwidth}{!}{%
\begin{tikzpicture}[
  block/.style={rectangle, draw, fill=blue!8, text width=3.5cm, text centered, minimum height=1.2cm, rounded corners=3pt, font=\small},
  greenblock/.style={rectangle, draw, fill=green!10, text width=3.5cm, text centered, minimum height=1.2cm, rounded corners=3pt, font=\small},
  arrow/.style={-{Stealth[length=3mm]}, thick}
]
\node[block] (china) at (0,0) {Chinese Smart Park Models\\[2pt]\scriptsize Core Features: OSS, Workflow, Communication};
\node[block] (adapt) at (6,0.8) {Technology Adaptation Framework};
\node[block] (offline) at (6,-0.8) {Offline-First Architecture};
\node[block] (ethiopia) at (12,0.8) {Ethiopian IPDC Context};
\node[greenblock] (platform) at (12,-0.8) {IPDC Digital OSS Platform};

\draw[arrow] (china) -- (adapt);
\draw[arrow] (china) -- (offline);
\draw[arrow] (adapt) -- (ethiopia);
\draw[arrow] (adapt) -- (platform);
\draw[arrow] (offline) -- (platform);
\draw[arrow] (ethiopia) -- (platform);
\end{tikzpicture}}
\caption{Conceptual Framework for Technology Adaptation}
\label{fig:conceptual-framework}
\end{figure}

Chinese smart park experience defines what digital park governance can look like in practice. The adaptation framework answers what can be transferred and what needs rethinking. Offline-first architecture solves the connectivity problem that makes direct replication impossible. The Ethiopian operating environment sets the requirements that all three must satisfy. The IPDC Digital OSS Platform is where those influences converge.

None of the reviewed literature addresses the full combination: complex industrial workflows, multiple stakeholder roles, and variable connectivity. That gap is where this research starts, and the methodology described in the next chapter was designed specifically to work within it.
